\rhead{MN-MED: Biomedical Signals \& Systems}
\phantomsection
\section*{Biomedical Signals \& Systems (MN-MED)}
\label{minor:MED_L1}

\noindent \small{\hyperref[main:scheme]{$\leftarrow$ Back to Scheme}} \vspace{0.5cm}

\noindent \textbf{Credits:} 2 (2-0-0)

\subsection*{Course Content}
\noindent\textbf{Unit 1} \\
\textit{The Body as a Data Source:} Biomedical signals as time-series data generated by physiological processes: electrical, mechanical, acoustic, and chemical origins; Classification of biomedical signals: deterministic vs.\ stochastic, continuous vs.\ discrete, stationary vs.\ non-stationary; Key signal types and their clinical context: ECG (cardiac electrical activity), EEG (brain waves), EMG (muscle activity), PPG (blood volume pulse), and PCG (heart sounds); Sampling and digitization: the Nyquist criterion applied to physiological bandwidths; Signal quality metrics: SNR, baseline wander, and motion artifacts as noise modeling problems; Introduction to standard biomedical data formats: EDF, HL7, and DICOM headers for time-series records.

\vspace{0.3cm}

\noindent\textbf{Unit 2} \\
\textit{Signal Processing Foundations for Physiological Data:} The Fourier Transform as the bridge between time-domain waveforms and frequency-domain physiology: interpreting spectral content of ECG and EEG; Short-Time Fourier Transform (STFT) and spectrograms for non-stationary signals; Discrete Wavelet Transform (DWT) for multi-resolution analysis: decomposing EEG into delta, theta, alpha, beta, and gamma bands; Digital filtering: FIR and IIR filter design for artifact removal (powerline interference at 50/60 Hz, baseline wander); Convolution as the computational primitive of linear filtering; Feature extraction from time-domain (mean, variance, zero-crossing rate) and frequency-domain (band power, spectral entropy) for downstream analysis.

\vspace{0.3cm}

\noindent\textbf{Unit 3} \\
\textit{Physiological Rhythms and Event Detection:} Cardiac cycle anatomy and the PQRST complex as a structured pattern in ECG; Pan-Tompkins algorithm for QRS detection: bandpass filtering, differentiation, squaring, and moving-window integration as a classical signal processing pipeline; R-R interval time series and Heart Rate Variability (HRV) as a window into the autonomic nervous system; Seizure detection in EEG: energy-based and threshold-based detectors; Apnea detection from respiratory signals using autocorrelation and spectral analysis; The general paradigm: segmentation $\rightarrow$ feature extraction $\rightarrow$ event classification as a reusable computational template across all biosignal types.

\vspace{0.3cm}

\noindent\textbf{Unit 4} \\
\textit{Biosensor Systems and Data Acquisition Pipelines:} Transducer principles: converting physiological quantities into electrical signals (electrodes, pressure transducers, optical sensors); Instrumentation amplifiers, differential amplification, and common-mode rejection as analog front-end design concepts; Analog-to-Digital conversion pipeline: anti-aliasing filters, ADC resolution, and dynamic range; Wearable sensor systems: IMUs, pulse oximeters, and continuous glucose monitors as constrained embedded data acquisition devices; Data streaming and buffering: circular buffers and real-time processing constraints; Introduction to IoT-based patient monitoring architectures: edge preprocessing, MQTT-based transmission, and cloud aggregation.

\vspace{0.3cm}

\noindent\textbf{Unit 5} \\
\textit{Clinical Decision Support and Signal Intelligence:} The clinical decision support system (CDSS) as a software pipeline: signal ingestion $\rightarrow$ preprocessing $\rightarrow$ feature engineering $\rightarrow$ classification $\rightarrow$ alert generation; Arrhythmia classification from ECG as a canonical supervised learning problem: feature-based classifiers vs.\ end-to-end approaches; Stress and fatigue detection from HRV and EEG features; Alarm fatigue in ICUs: the false-positive problem and precision-recall tradeoffs in safety-critical systems; Introduction to interoperability standards: FHIR as a RESTful API standard for exchanging clinical data; Regulatory and ethical constraints: FDA classification of Software as a Medical Device (SaMD) and data privacy under HIPAA.

\newpage
\rhead{Curriculum Schema}
